% Part: history
% Chapter: biographies
% Section: gerhard-gentzen

\documentclass[../../../include/open-logic-section]{subfiles}

\begin{document}

\olfileid[hi]{his}{bio}{gen}

\olsection{गेरहार्ड गेंट्सेन}

\olphoto{gentzen-gerhard}{गेरहार्ड गेंट्सेन}

गेरहार्ड गेंट्सेन मुख्यतः संरचनात्मक प्रमाण सिद्धांत के रचयिता, और विशेष
रूप से प्राकृतिक निगमन तथा अनुक्रम कलन की !!{derivation} प्रणालियों के
निर्माता के रूप में जाने जाते हैं। उनका जन्म 24 नवंबर 1909 को जर्मनी के
ग्राइफ़्सवाल्ड में हुआ। प्रारंभिक विद्यालय जाने से पहले गेरहार्ड ने तीन वर्ष
घर पर शिक्षा पाई; विद्यालय में वे शिक्षा की दृष्टि से अधिकतर सहपाठियों से
पीछे थे। इसके बावजूद वे प्रतिभाशाली विद्यार्थी थे और उन्होंने गणित की प्रबल
योग्यता दिखाई। उनकी रुचियाँ विविध थीं; उदाहरणतः वे अपनी माँ के लिए कविताएँ
और विद्यालय के रंगमंच के लिए नाटक भी लिखते थे।

गेंट्सेन ने विश्वविद्यालयी अध्ययन ग्राइफ़्सवाल्ड विश्वविद्यालय में आरंभ
किया, किंतु फिर गॉटिंगेन, म्यूनिख और बर्लिन जाते रहे। उन्होंने 1933 में हरमन
वाइल के निर्देशन में गॉटिंगेन विश्वविद्यालय से डॉक्टरेट प्राप्त की। (पॉल
बर्नाइस ने उनके अधिकांश कार्य का पर्यवेक्षण किया था, किंतु नाज़ियों ने उन्हें
विश्वविद्यालय से निकाल दिया।) 1934 में गेंट्सेन ने डेविड हिल्बर्ट के सहायक
के रूप में काम आरंभ किया। उसी वर्ष अपने शोधपत्रों \emph{Untersuchungen
\"{u}ber das logische Schlie\ss en I--II [Investigations Into Logical
Deduction I--II]} में उन्होंने अनुक्रम कलन और प्राकृतिक निगमन की
!!{derivation} प्रणालियाँ विकसित कीं। 1936 में उन्होंने पियानो अभिगृहीतों की
संगतता सिद्ध की।

नाज़ियों के साथ गेंट्सेन का संबंध जटिल है। जिस समय उनके मार्गदर्शक बर्नाइस
को जर्मनी छोड़ने के लिए विवश किया गया, उसी समय गेंट्सेन नाज़ी अर्धसैनिक
संगठन एसए की विश्वविद्यालय शाखा में सम्मिलित हुए। अनेक जर्मनों की तरह वे
नाज़ी दल के सदस्य थे। युद्ध के दौरान उन्होंने वायु गुप्तचर इकाई में दूरसंचार
अधिकारी के रूप में सेवा की। किंतु 1942 में मानसिक विक्षोभ के कारण उन्हें
सेवा से मुक्त कर दिया गया। यह स्पष्ट नहीं कि गेंट्सेन की निष्ठा नाज़ी दल के
साथ थी अथवा उन्होंने अकादमिक सफलता सुनिश्चित करने के लिए दल की सदस्यता ली
थी।

1943 में गेंट्सेन को जर्मन यूनिवर्सिटी ऑफ प्राग के गणितीय संस्थान में
अकादमिक पद का प्रस्ताव मिला, जिसे उन्होंने स्वीकार किया। किंतु 1945 में
प्राग के नागरिकों ने जर्मन अधिग्रहण के विरुद्ध विद्रोह किया। सोवियत सेनाएँ
नगर पहुँचीं और विश्वविद्यालय के सभी प्राध्यापकों को गिरफ़्तार कर लिया। नाज़ी
संगठनों की सदस्यता के कारण गेंट्सेन को बेगार शिविर ले जाया गया। 4 अगस्त 1945
को 35 वर्ष की आयु में अपनी कोठरी में कुपोषण से उनकी मृत्यु हुई।

\begin{reading}
गेंट्सेन की पूर्ण जीवनी के लिए \citet{Menzler-Trott2007} देखें। नाज़ी शासन
के अधीन गणितज्ञों पर एक रोचक पुस्तक, जिसमें गेंट्सेन के जीवन का संक्षिप्त
विवरण भी है, \citet{Segal2014} ने दी है। तार्किक निगमन पर गेंट्सेन के
शोधपत्र मूल जर्मन में उपलब्ध हैं \citep{Gentzen1935a,Gentzen1935b}।
\citet{Gentzen1969} ने गेंट्सेन के शोधपत्रों के अंग्रेज़ी अनुवाद एक खंड में
संकलित किए हैं; इसमें एक जीवनी-रेखा भी है।
\end{reading}

\end{document}
